\section{Introduction}

\label{sec:intro}
%% ============================================================================

Building production infrastructure for AI-native blockchain applications requires solving problems across at least sixteen distinct technology domains: identity and access management, API gateway routing, application runtime, secrets management, object storage, vector search, analytics, payment processing, distributed consensus, post-quantum cryptography, multi-party computation, fully homomorphic encryption, GPU-accelerated computing, wire-level network protocols, AI model inference, and AI agent orchestration.

The conventional approach assembles point solutions: Auth0 for identity, Kong for API routing, Supabase for the backend, HashiCorp Vault for secrets, Pinecone for vector search, and so on. Each point solution has its own authentication model, deployment topology, monitoring surface, and upgrade cadence. The integration cost between $n$ services scales as $O(n^2)$ in the worst case, and the operational burden grows with each additional vendor dependency.

The Hanzo + Lux stack takes the opposite approach. Every layer---from the identity provider to the consensus engine to the FHE runtime---is designed as a single integrated system. Services share a common identity model (Hanzo IAM OIDC tokens), a common secrets backend (Hanzo KMS), a common observability pipeline (Prometheus + Victoria Metrics), and a common deployment mechanism (Kubernetes operators with 18 custom resource definitions).

\paragraph{Scope.} This paper compares each layer of the Hanzo + Lux stack against the most common alternatives as of Q1~2026. We evaluate on five axes: (1)~feature completeness, (2)~raw performance, (3)~integration cost measured in lines of glue code and setup time, (4)~total cost of ownership at scale, and (5)~unique differentiators that cannot be replicated by combining alternatives.

\paragraph{Non-goals.} We do not claim that every Hanzo component is the best possible solution in isolation. A team building only a CRUD application with no blockchain or AI requirements would be well-served by Supabase or Firebase. The advantage of the Hanzo + Lux stack emerges when the application requires \textit{multiple} categories simultaneously---which is precisely the case for AI-native blockchain infrastructure.

\paragraph{Organization.} Sections~\ref{sec:identity}--\ref{sec:agents} each address one technology category with a feature matrix, performance data, integration cost analysis, and verdict. Section~\ref{sec:integrated} quantifies the compounding advantage of integration. Section~\ref{sec:deployment} describes the unified deployment model. Section~\ref{sec:conclusion} summarizes findings and provides a decision framework.

%% ============================================================================
